\subsection{Baseline tracking with fp32 input and output}
\label{app:fp32_baseline_tracking}
Table~\ref{tab:fp32_baseline_tracking} compares the three baseline recoveries
on H100 PCIe with fp32 input and output. It uses copied weights, output checks
before timing, three-round medians and the first comparison-group attempt
passing the validity and noise gates. Kernel timings use the aligned scope;
mixer timings cover an eager forward and backward of the full module.
The SSD recoveries here use $R=O=1$; fp32 memory is not measured.
HM timings retain the earlier backward-buffer schedule; the buffer
optimization was measured only in bf16.

\begin{table}[!htbp]
\centering\scriptsize
\caption{Baseline tracking with fp32 input/output: time ratios ours/official
($<1$ means ours is faster), geometric mean with min--max underneath over four
Large or six LM shapes. Mamba~1 uses IEEE arithmetic, Mamba~2 uses tf32,
and our Mamba~3 recovery uses its opt-in bf16-operand mode to match the
official implementation's internal operand casts. The Mamba~3 row is not
a comparison of full-fp32 scan arithmetic. PyTorch~2.14, Triton~3.8.}
\label{tab:fp32_baseline_tracking}
\adjustbox{max width=\linewidth}{% Generated directly from verified anonymous-revision-14 JSONL by appendix_fp32_baseline/build.py.
\begin{tabular}{@{}l ccc ccc@{}}
\toprule
 & \multicolumn{3}{c}{Large} & \multicolumn{3}{c}{LM} \\
\cmidrule(lr){2-4}\cmidrule(lr){5-7}
Pair (our operands) & kernel fwd & kernel f+b & mixer f+b & kernel fwd & kernel f+b & mixer f+b \\
\midrule
Mamba 1 (IEEE) & \shortstack{0.54\\{\tiny (0.51--0.65)}} & \shortstack{0.53\\{\tiny (0.45--0.65)}} & \shortstack{0.97\\{\tiny (0.87--1.02)}} & \shortstack{0.58\\{\tiny (0.51--0.67)}} & \shortstack{0.61\\{\tiny (0.50--0.68)}} & \shortstack{0.99\\{\tiny (0.95--1.03)}} \\
Mamba 2 (tf32) & \shortstack{1.36\\{\tiny (1.16--1.58)}} & \shortstack{1.00\\{\tiny (0.97--1.08)}} & \shortstack{1.15\\{\tiny (1.11--1.19)}} & \shortstack{1.54\\{\tiny (1.15--1.70)}} & \shortstack{1.02\\{\tiny (0.92--1.09)}} & \shortstack{1.18\\{\tiny (1.15--1.22)}} \\
Mamba 3 (bf16 operands) & \shortstack{0.88\\{\tiny (0.76--0.98)}} & \shortstack{0.94\\{\tiny (0.85--1.02)}} & \shortstack{0.95\\{\tiny (0.91--0.98)}} & \shortstack{0.88\\{\tiny (0.76--1.11)}} & \shortstack{0.94\\{\tiny (0.85--1.12)}} & \shortstack{0.94\\{\tiny (0.85--1.02)}} \\
\bottomrule
\end{tabular}
}
\end{table}

With our default tf32 operands instead, the Mamba~3 kernel forward-and-backward
ratio is $4.03$ on Large and $3.57$ on LM, and the eager mixer ratio is $2.00$
and $2.06$, respectively. These ratios describe the measured operand configurations; no subsequent
fp32 tuning is included.
